% fluids_capfbd: the forces that hold up a capillary column (Vol II Ch 2, fig:fluids-capfbd)
% Left: water risen to height h in a tube of radius r; the box marks the contact line.
% Right (zoom of the left wall): the liquid surface leaves the glass at the contact
% angle theta = 30 deg, measured from the wall (down) to the surface's tangent.
% The surface pulls the glass along its tangent, S per unit length (dashed, down and
% away from the wall).  By the third law the glass pulls the liquid with S the other
% way: S cos(theta) up the wall and S sin(theta) towards the wall.  Round the circular
% rim the horizontal parts cancel; the vertical parts give 2 pi r S cos(theta) up,
% which holds up the weight rho g pi r^2 h of the column.
\begin{tikzpicture}[x=1cm,y=1cm,
    water/.style={fill=cPrim!12,draw=none},
    glass/.style={fill=cMoon,draw=cSteel,thin line}]
  % ---------------- the tube ----------------
  \begin{scope}
    \def\r{0.32}\def\w{0.09}\def\h{2.0}
    \fill[water] (-1.4,-1.0) rectangle (1.4,0);
    \fill[water] (-\r,0) -- (-\r,\h+0.22) .. controls (-0.12,\h) and (0.12,\h) .. (\r,\h+0.22) -- (\r,0) -- cycle;
    \draw[cPrim,thin line] (-1.4,0) -- ({-\r-\w},0) ({\r+\w},0) -- (1.4,0);
    \draw[cPrim,thin line] (-\r,\h+0.22) .. controls (-0.12,\h) and (0.12,\h) .. (\r,\h+0.22);
    \draw[glass] ({-\r-\w},-0.7) rectangle (-\r,2.75);
    \draw[glass] (\r,-0.7) rectangle ({\r+\w},2.75);
    \draw[outline] (-1.4,0.3) -- (-1.4,-1.0) -- (1.4,-1.0) -- (1.4,0.3);
    \draw[dim] (-0.95,0) -- (-0.95,\h) node[midway,anchor=east,cSteel] {$h$};
    \draw[dimtick] ({-\r-\w-0.05},\h) -- (-1.1,\h);
    % zoom box round the right-hand contact line
    \draw[cGrey,thin line,dash pattern=on 2pt off 1.5pt] (0.05,\h-0.12) rectangle (0.62,\h+0.52);
    \draw[cGrey,thin line,dash pattern=on 2pt off 1.5pt] (0.62,\h+0.52) -- (1.75,2.95);
    \draw[cGrey,thin line,dash pattern=on 2pt off 1.5pt] (0.62,\h-0.12) -- (1.75,-0.25);
  \end{scope}
  % ---------------- zoom on the contact line (right-hand wall) ----------------
  \begin{scope}[xshift=5.0cm,yshift=0.1cm]
    \def\th{30}
    \coordinate (C) at (0,1.6);
    \fill[water] (C) .. controls ($(C)+({-0.6*sin(\th)},{-0.6*cos(\th)})$) and (-1.3,0.5) .. (-3.1,0.5)
      -- (-3.1,-0.3) -- (0,-0.3) -- cycle;
    \draw[cPrim,construction] (C) .. controls ($(C)+({-0.6*sin(\th)},{-0.6*cos(\th)})$) and (-1.3,0.5) .. (-3.1,0.5);
    \draw[glass] (0,-0.3) rectangle (0.5,2.9);
    \node[anchor=south,cSteel] at (0.25,2.9) {glass};
    \node[cPrim] at (-2.3,0.1) {liquid};
    % angle of contact
    \draw[angle arc] ($(C)+(0,-0.55)$) arc[start angle=-90,end angle={-90-\th},radius=0.55];
    \node at ($(C)+({-0.78*sin(\th/2)},{-0.78*cos(\th/2)})$) {$\theta$};
    % pull of the surface on the glass (along the tangent)
    \draw[force,dash pattern=on 3pt off 1.5pt] (C) -- ($(C)+({-1.2*sin(\th)},{-1.2*cos(\th)})$);
    \node[anchor=north,cCrim,align=center,font=\footnotesize] at ($(C)+({-1.2*sin(\th)},{-1.25*cos(\th)})$) {$S$ on the glass};
    % pull of the glass on the liquid (third-law partner) and its components
    \draw[component,cCrim] (C) -- ($(C)+(0,{1.2*cos(\th)})$);
    \node[anchor=east,cCrim] at ($(C)+(-0.05,{0.7*cos(\th)})$) {$S\cos\theta$};
    \draw[component,cCrim] (C) -- ($(C)+({1.2*sin(\th)},0)$);
    \node[anchor=west,cCrim,knockout] at ($(C)+({1.2*sin(\th)+0.05},-0.02)$) {$S\sin\theta$};
    \draw[force] (C) -- ($(C)+({1.2*sin(\th)},{1.2*cos(\th)})$) node[anchor=south west,inner sep=1pt] {$S$};
    \node[contact point] at (C) {};
  \end{scope}
\end{tikzpicture}
